
% =====================================================================
\section{Scope and conclusion}
\label{sec:scope}
% =====================================================================

\subsection{What has been proved}

From the single hypothesis of a bounded resolvable space (\Cref{ax:brs}) we have
established the following, each from stated definitions:

\begin{enumerate}[label=(\arabic*),leftmargin=2.4em]
\item \textbf{Individuation by negation} (\Cref{thm:negation-unique}): in a whole
with no distinguished element, a part can be fixed only as the complement of its
negation-sequence; this is the unique selector-free individuation.
\item \textbf{Non-instantaneity and residue} (\Cref{thm:noninstant}): no
individuation of a genuine part completes in zero steps, and at every stage before
completion a separator of measure $\ge\thick$ remains uncommitted.
\item \textbf{Boundary thickness} (\Cref{thm:thickness}): every realisable
separator has measure $\ge\mu_{\min}>0$; no sharp cut exists for an embedded
agent. The floor $\thick$ is strictly positive.
\item \textbf{Three agreeing derivations} (\Cref{thm:floor-agree}): the
geometric, representational, and cost-limited floors are all positive and bound
the same constant.
\item \textbf{Detectability and the sorites} (\Cref{thm:detectability},
\Cref{cor:sorites}): a region is distinguishable iff its scale exceeds its
boundary thickness; the sorites dissolves as emergence across a band of width
$\thick$, with no decisive grain.
\item \textbf{Locating versus visiting} (\Cref{thm:asymmetry},
\Cref{thm:locating-cost}): individuating a region is costly and the cost grows
with the logarithm of the region's rarity, while a measure-preserving traversal
visits the region freely and scales with rarity only in frequency; three invalid
inferences transferring the former onto the latter are classified and blocked
(\Cref{prop:one-form}).
\item \textbf{Non-self-verifiability} (\Cref{thm:diagonal},
\Cref{thm:residue-not-measure}): no verifier is total and closed under its own
diagonal, an unconditional self-contained result; and \emph{under the explicit
closure hypothesis} \ref{sa:closed}---that a behaviour-defined collection of
regions is again a region---the irreducible residue is not an exhibitable set and
the floor can only be bounded below, never displayed. This last upgrade is stated
conditionally; the positivity of the floor (items 3--4) does not rely on it.
\item \textbf{Non-return and cessation} (\Cref{thm:nonreturn},
\Cref{thm:cessation}): the committed record is monotone, so no process returns to
a prior individuation state; thinning a separator has diverging cost and bounded
return, so a finite agent halts at the floor and knows only bounded wholes.
\end{enumerate}

These results turn on one constant, the floor $\thick>0$. Its positivity is
reached from four independent directions---topological, representational,
cardinality-theoretic, and economic---each shown to bound or price the same
quantity unconditionally. Two further directions speak to its character rather
than its positivity: the self-referential, which under the closure hypothesis
\ref{sa:closed} shows the floor cannot be displayed as a located set, and the
order-theoretic, which shows the record of individuations cannot return.

\subsection{Scope}

The theorems are statements about an abstract bounded resolvable space and an
embedded agent that individuates parts within it by negation. We have made no
claim about any particular system, and the proofs use only the metric--measure
and order structure declared in \Cref{sec:setting}. The hypotheses
(\Cref{ax:brs}) are deliberately minimal: boundedness with compact closure,
finite resolution, and hierarchical refinability. Where each clause is used has
been marked, and the necessity of boundedness is shown by a counterexample
(\Cref{prop:bdd-necessary}). Material that could not be proved at this standard
has been omitted rather than asserted.

One hypothesis beyond \Cref{ax:brs} is used, and only in \Cref{sec:residue}: the
closure hypothesis \ref{sa:closed}, that a collection of regions defined by a
verifier's own behaviour is again a region of $\Space$. We do not prove
\ref{sa:closed}; we isolate it, argue that an embedding makes it natural, note
that discharging it in general would require exhibiting the measurable encoding
and point-surjection that the classical fixed-point theorem of~\citet{lawvere1969}
needs, and mark every conclusion that depends on it. The positive floor and all
results of \Cref{sec:thickness,sec:floor,sec:detect,sec:locating,sec:inquiry}
stand independently of \ref{sa:closed}.

\subsection{Conclusion}

A bounded whole, finitely resolved, cannot be divided into a part and its
complement without depositing positive content into the separator, and that
content cannot be driven to zero---and, where the agent's own verifying behaviour
is itself a region, cannot be certified from within. The part is
individuated by what it is not; the separation is never instantaneous and never
exact; the residue is irreducible; the record of separations only grows; and a
finite agent rationally halts at a positive floor, knowing bounded wholes and not
their sub-floor parts. The floor $\thick>0$ is the single invariant through which
individuation, boundary thickness, detectability, verification, and non-return
all speak. It is the irreducible boundary of a bounded resolvable space.

% =====================================================================
\bibliographystyle{plainnat}
\bibliography{references}
% =====================================================================

\end{document}
